% Two-page project report for the DWM practical.
% Build: tectonic DWM_Project_Report.tex
\documentclass[11pt,a4paper]{article}
\usepackage[T1]{fontenc}
\usepackage{newtxtext,newtxmath}
\usepackage[a4paper,margin=1.8cm,top=1.6cm,bottom=1.6cm]{geometry}
\usepackage{booktabs}
\usepackage{enumitem}
\usepackage{titlesec}
\usepackage[hidelinks]{hyperref}
\setlist{nosep,leftmargin=1.2em}
\titleformat{\section}{\large\bfseries}{\thesection.}{0.5em}{}
\titlespacing*{\section}{0pt}{9pt}{4pt}
\setlength{\parskip}{4pt}
\setlength{\parindent}{0pt}
\pagestyle{empty}

\begin{document}

\begin{center}
{\Large\bfseries Sentinel: AI-Based Intrusion Detection with Human Review and a Data Warehouse}\\[5pt]
Project Report -- Data Warehousing and Mining (DWM) Practical\\[3pt]
Jai Chavan \quad Jeet Mane \quad Amaan Khan\\[1pt]
Guide: Dr. Kavita Shirsat \quad | \quad Department of Computer Engineering, Vidyalankar Institute of Technology, Mumbai\\[1pt]
Project link: \url{https://github.com/JaiChavan040906/IDS}
\end{center}

\section{Abstract}
Sentinel is a web application that detects attacks in network traffic. A CNN-LSTM model
classifies every network flow, and flows that the model is not sure about are sent to a security
analyst instead of being decided automatically. The analyst's decisions are saved and can be used
to retrain the model after an administrator approves them. All traffic, alerts and reviews are
also loaded into a data warehouse, where they can be analysed with OLAP operations. On unseen
test data, only 1.48\% of flows needed the analyst, and the automatic decisions were 99.999\%
correct.

\section{Problem Statement and Objectives}
AI-based intrusion detection systems give a label to every flow, even for attacks they have never
seen, and a wrong label can look just as confident as a correct one. The alerts are also stored in
an ordinary database, which makes it hard to analyse them by day, attack type or service. The
objectives of this project are:
\begin{itemize}
\item to detect normal traffic and 12 types of network attacks using a CNN-LSTM model;
\item to send only uncertain flows to a security analyst and save the analyst's decisions;
\item to improve the model from verified feedback, only after admin approval;
\item to build a data warehouse with ETL and analyse the data using OLAP operations.
\end{itemize}

\section{Methodology}
\textbf{1. Dataset.} We used the Improved CIC-IDS2017 dataset (five days of lab traffic). After
removing duplicate rows and columns such as IP addresses and ports, 1,714,812 flows with 64
features and 13 classes (normal traffic and 12 attack types) remained. The data was split into
70\% training, 15\% validation and 15\% test. The test part was kept aside and used only once,
for the final evaluation.

\textbf{2. Detection models.} The main model is a CNN-LSTM network. The 64 features of a flow are
read as a sequence and passed through a 1-D convolution (64 filters), max-pooling, a second
convolution (128 filters), an LSTM layer (128 units), dropout (0.3) and an output layer with one
value per class. It was trained with the AdamW optimiser (learning rate 0.001, batch size 1,024)
for up to 20 epochs, stopping early when the validation score stopped improving. Classes with
fewer than 1,000 training rows were oversampled so the model did not ignore them. A Random Forest
(100 trees) and LightGBM were trained on the same data for comparison, and the Random Forest is
also used as a second opinion. Every model was trained with three random seeds.

\textbf{3. Uncertainty check.} A flow is sent to the analyst if the prediction is not confident
(high entropy), if the flow is unlike the training data (Mahalanobis distance), or if the Random
Forest gives a different answer. All other flows are closed or alerted automatically.

\textbf{4. Analyst review and retraining.} The analyst marks a flow as malicious, benign or
unsure and gives a reason. These cases are stored and shown as similar incidents later. After
20 decisions are verified by the admin, a new model can be trained; it is deployed only if it
passes fixed checks, and it can be rolled back.

\textbf{5. Data warehouse and ETL.} The ETL job \emph{extracts} the dataset and the portal's alerts
and reviews, \emph{transforms} them (time keys, port to service name, attack to category) and
\emph{loads} them into a separate SQLite warehouse. The main design is a star schema: the fact
table \texttt{fact\_flow} (one row per flow) with time, attack, service, protocol and split
dimensions. Snowflake and fact constellation versions, an hourly aggregate table and two data marts
are also built. A full load rebuilds everything; an incremental load adds only new alerts.

\textbf{6. OLAP.} The OLAP page supports roll-up, drill-down, slice, dice and pivot, and shows the
result as a table, a chart and a 3-D cube.

\section{Results}
All detection results were measured on the 15\% test data (257,222 flows), which the models never
saw during training. Table 1 compares the models (attack vs normal, average of three runs).

\noindent\begin{minipage}{\linewidth}
\begin{center}
\textbf{Table 1.} Model comparison on the test data\\[3pt]
\begin{tabular}{@{}lcccc@{}}
\toprule
Model & Detection rate & Precision & F1 score & Accuracy (13 classes) \\
\midrule
Reference paper (de Nascimento and Hou, 2025) & 96.6\% & 92.1\% & 94.3\% & -- \\
Random Forest & 99.92\% & 99.90\% & 99.91\% & 99.97\% \\
LightGBM & 99.91\% & 99.84\% & 99.88\% & 99.96\% \\
CNN-LSTM (our main model) & 99.75\% & 99.71\% & 99.73\% & 99.89\% \\
CNN only (LSTM removed) & 99.23\% & 99.70\% & 99.46\% & 99.80\% \\
\bottomrule
\end{tabular}
\end{center}
\end{minipage}

All our models scored above the reference paper. The Random Forest was slightly better than the
CNN-LSTM, and removing the LSTM layer lowered the F1 score, so the LSTM does help. Table 2 shows
how the full system (model plus uncertainty check) handled the test traffic, and the warehouse
measurements.

\begin{center}
\textbf{Table 2.} System and warehouse results\\[3pt]
\begin{tabular}{@{}lr@{\hspace{3em}}lr@{}}
\toprule
\multicolumn{2}{@{}l}{\textbf{Detection}} & \multicolumn{2}{l@{}}{\textbf{Data warehouse}} \\
\midrule
Flows sent to analyst & 1.48\% & Flows in fact table & 1,714,812 \\
Model errors caught in review & 98.9\% & Full ETL load & about 25 s \\
Accuracy of automatic decisions & 99.999\% & Incremental load & about 0.3 s \\
Attacks missed / false alarms & 2 / 0 & Query on aggregate vs all flows & 0.3 ms vs 2.3 s \\
\bottomrule
\end{tabular}
\end{center}

Table 3 shows the OLAP operations performed on the flow cube of the warehouse and what each one
showed about the traffic.

\begin{center}
\textbf{Table 3.} OLAP operations on the flow cube (measure: number of flows)\\[3pt]
\begin{tabular}{@{}lll@{}}
\toprule
Operation & What was done & Result \\
\midrule
Roll-up & Attack class $\to$ category, per day & Monday has no attacks; Wednesday mostly DoS (177,506) \\
Drill-down & Wednesday: day $\to$ hour & 158,662 DoS flows fall in one hour \\
Slice & Service = HTTP & 94.4\% of all attack flows used HTTP \\
Dice & Wed--Thu, DoS/DDoS, HTTP/HTTPS & a smaller sub-cube of web flooding attacks \\
Pivot & Swap rows and columns & same numbers seen from the other dimension \\
\bottomrule
\end{tabular}
\end{center}

\textbf{Tools used:} Python, PyTorch, scikit-learn, LightGBM, pandas, FastAPI, SQLite, React and
Recharts. The full source code is available at the project link above.

\section{Conclusion and Future Work}
The system handles most traffic on its own and asks a human only about doubtful flows, which
keeps the automatic decisions almost error-free. The data warehouse makes the traffic and alert
history easy to explore with OLAP. In future, we plan to detect slow attacks such as SSH brute
force by using features over many flows, and to test the system on the CSE-CIC-IDS2018 dataset.

\end{document}
